% Generated draft: requires mathematical and source review.
\documentclass[11pt]{article}
\usepackage[margin=25mm]{geometry}
\usepackage{fontspec}
\setmainfont{Latin Modern Roman}
\usepackage{amsmath,amssymb,mathtools,unicode-math}
\setmathfont{Latin Modern Math}
\usepackage{enumitem,needspace,xurl,hyperref}
\setlength{\parindent}{0pt}
\setlength{\parskip}{5pt}
\setlength{\emergencystretch}{3em}
\newcommand{\N}{\mathbb{N}}
\newcommand{\Z}{\mathbb{Z}}
\newcommand{\Q}{\mathbb{Q}}
\newcommand{\R}{\mathbb{R}}
\newcommand{\C}{\mathbb{C}}
\newcommand{\F}{\mathbb{F}}
\newcommand{\A}{\mathbb{A}}
\newcommand{\PP}{\mathbb{P}}
\newcommand{\E}{\mathbb{E}}
\newcommand{\eps}{\varepsilon}
\newcommand{\dd}{\,\mathrm{d}}
\begin{document}
\section*{UnsolvedMath ID: 30000108}
\textbf{Source problem code:} catalogue entry
\subsection{Definitions and mathematical statement}
Let \(\mathfrak g\) be a complex simple Lie algebra of one of the classical types \(A_r,B_r,C_r,D_r\), with a valid rank for that type. For a dominant integral weight \(\eta\), let \(V_\eta\) be the irreducible finite-dimensional highest-weight representation. Define the Clebsch--Gordan multiplicity by
\[
C_{\lambda,\mu}^{\nu}=\dim_{\C}\operatorname{Hom}_{\mathfrak g}(V_\nu,V_\lambda\otimes V_\mu).
\]
For fixed dominant integral weights \(\lambda,\mu,\nu\), consider the stretching function
\[
q(N)=C_{N\lambda,N\mu}^{N\nu},\qquad N\in\N,\quad N\geq1.
\]
The source describes it as an Ehrhart quasipolynomial. If \(m\geq1\) is a period, its constituent polynomials \(Q_0(T),\ldots,Q_{m-1}(T)\in\Q[T]\) satisfy
\[
q(N)=Q_j(N)\quad\text{whenever }N\geq1\text{ and }N\equiv j\pmod m.
\]
The source conjecture is that, for every indicated Lie type and rank and every triple of dominant integral weights, each constituent has nonnegative coefficients in the ordinary monomial basis:
\[
Q_j(T)=\sum_{\ell\geq0}a_{j,\ell}T^\ell,
\qquad a_{j,\ell}\geq0\quad(0\leq j<m,\ \ell\geq0).
\]
Identically zero constituents are allowed. This is stronger than the automatic nonnegativity of representation multiplicities at integer arguments. The variable is the original stretching factor \(N\), not the reparameterized residue variable in \(N=j+mk\); coefficient positivity after that change is a different statement. The formulation uses positive stretches and makes no extra claim identifying constituent values at zero with \(C_{0,0}^{0}\).

The supplied assessment checked 2026-08-21 reports that Rassart proves polynomiality in the relevant chambers and that later work treats restricted cases, while general coefficient nonnegativity remains unproved in that check. It does not reproduce the exact chamber, type, or parameter hypotheses of those results; it does not justify asserting a single polynomial in every classical type.
\subsection{Short English statement}
When all three highest weights are multiplied by \(N\), the resulting tensor-product multiplicity is described by constituent polynomials on residue classes. The conjecture asks that every coefficient of each polynomial, in powers of \(N\), be nonnegative.
\subsection{Sources}
\begin{itemize}
\item Local UnsolvedMath registry, permanent ID 30000108, catalogue entry, \emph{Nonnegative Stretched Clebsch--Gordan Coefficients}.
\item \emph{Ehrhart Quasipolynomials: Algebra, Combinatorics, and Geometry} (2005), Oberwolfach Reports, pp. 2071--2102, DOI \url{https://doi.org/10.4171/owr/2004/39}.
\item Rassart, supplied preprint link \url{https://arxiv.org/abs/math/0308101}.
\item King, supplied article link \url{https://www.numdam.org/item/ALCO_2022__5_6_1227_0.pdf}.
\item Supplied literature triage checked 2026-08-21 and statement-verification claim of comparison with the official report; exact theorem numbers and full article titles are not supplied.
\end{itemize}
\end{document}
